\section{Transformers y mecanismos de atención}
\label{sec:marco_atencion}

Para las secuencias, como un texto, las primeras redes avanzaban un elemento por vez, de modo que el acceso a un
elemento lejano pasaba por todos los intermedios y el cálculo de un paso esperaba al del anterior. En 2017, la
atención resolvió ese límite, y la arquitectura construida sobre ella, el transformer, es hoy la base de los
modelos de texto~\cite{vaswani2017attention} y de buena parte de los de imagen~\cite{dosovitskiy2021image}.

\subsection{Atención escalada por producto punto}
\label{sec:marco_atencion_producto}

La atención compara directamente todas las posiciones de una secuencia: cada posición de la salida pondera
los valores de todas las posiciones de la entrada según la compatibilidad entre una consulta ($Q$) y una clave
($K$). Los elementos de la secuencia se llaman \emph{tokens}, y son fragmentos de palabra en un texto o parches
en una imagen. Cada \emph{token} se representa con un \emph{embedding}, un vector aprendido, de modo que los
elementos parecidos quedan cerca en ese espacio. Como todas las posiciones se comparan a la vez, el cálculo se
paraleliza y aprovecha mejor el hardware que el recorrido secuencial~\cite{vaswani2017attention}.

La atención escalada por producto punto proyecta el embedding de cada token en tres espacios,
consulta $Q$, clave $K$ y valor $V$, y devuelve una combinación de valores pesada por la compatibilidad entre
consulta y clave:
\begin{equation}
    \mathrm{Attention}(Q, K, V) = \mathrm{softmax}\!\left(\frac{QK^{\top}}{\sqrt{d_k}}\right) V,
    \label{eq:marco_atencion}
\end{equation}
donde $d_k$ es la dimensión de la clave. La división por su raíz contrarresta que el producto punto
crezca con la dimensión y lleve la softmax a zonas de gradiente muy chico~\cite{vaswani2017attention}. En la Figura~\ref{fig:marco_atencion}, cada
fila de la matriz compara una consulta con las claves de todas las posiciones, y la softmax convierte esas
compatibilidades en pesos que suman uno.

% Figura: la matriz de pesos de atención.
% Se usa en sec:marco_atencion, después de la Ecuación de la atención escalada por producto punto.
% Los pesos son ilustrativos: muestran la forma de la operación, no un modelo entrenado.

\begin{figure}[H]
    \centering
    \begin{tikzpicture}[x=1mm, y=1mm]

        % Consultas (filas) y claves (columnas).
        \foreach \i in {1, 2, 3, 4} {
            \node[fignodo, minimum height=6mm, minimum width=9mm, inner sep=1pt,
                  draw=figazul!80!black, fill=figazul!12] (q\i) at (0, 30 - \i*8) {\footnotesize $q_{\i}$};
            \node[fignodo, minimum height=6mm, minimum width=6.5mm, inner sep=1pt,
                  draw=fignaranja!70!black, fill=fignaranja!18] (k\i) at (8 + \i*8, 34) {\footnotesize $k_{\i}$};
        }

        % Pesos softmax(QK^T / sqrt(d_k)); cada fila suma uno.
        \foreach \row/\vals in {1/{60,20,10,10}, 2/{15,55,20,10}, 3/{10,25,50,15}, 4/{5,10,25,60}} {
            \foreach \v [count=\col] in \vals {
                \fill[figceleste!\v!white] (4 + \col*8, 26 - \row*8) rectangle (12 + \col*8, 34 - \row*8);
                \node[font=\scriptsize] at (8 + \col*8, 30 - \row*8) {0.\ifnum\v<10 0\fi\v};
            }
        }
        \draw[figtrazo, shift={(12, -6)}] (0, 0) grid[step=8] (32, 32);

        % Salida.
        \foreach \i in {1, 2, 3, 4} {
            \node[fignodo, minimum height=6mm, minimum width=9mm, inner sep=1pt,
                  draw=figverde!70!black, fill=figverde!10, densely dashed] (o\i) at (66, 30 - \i*8)
                  {\footnotesize $o_{\i}$};
            \draw[figflecha] (45, 30 - \i*8) -- (o\i.west);
        }

        \node[font=\footnotesize, text=figtrazo] at (66, 32) {$o_i = \sum_j a_{ij} v_j$};
        \node[font=\footnotesize, text=figtrazo, align=center] at (28, -12)
            {pesos $a_{ij}$ = $\mathrm{softmax}\!\left(q_i k_j^{\top} / \sqrt{d_k}\right)$};

    \end{tikzpicture}
    \caption{Pesos de atención entre las posiciones de una secuencia de cuatro tokens (valores ilustrativos).}
    \label{fig:marco_atencion}
\end{figure}


La atención multi-cabeza proyecta $Q$, $K$ y $V$ en varios subespacios de menor dimensión, aplica
la Ecuación~\ref{eq:marco_atencion} en cada uno por separado y concatena los resultados, de modo que cada
cabeza captura relaciones de distinto tipo entre posiciones~\cite{vaswani2017attention}.

\subsection{Codificación de posición}
\label{sec:marco_posicion}

La Ecuación~\ref{eq:marco_atencion} no tiene en cuenta el orden de los tokens: si se permutan las posiciones de
la entrada, la salida se permuta de la misma manera y cada token recibe el mismo resultado. Para que el modelo
distinga «la lesión crece» de «crece la lesión», se suma al embedding de cada token un vector que depende de su
posición. El transformer original usa vectores fijos, formados por senos y cosenos de distinta frecuencia,
que dan a cada posición un vector único~\cite{vaswani2017attention}. La alternativa es aprender un vector por
posición, como un parámetro más, que es lo que hace el \gls{vit} con la posición de cada parche en la imagen
(Sección~\ref{sec:marco_vit})~\cite{dosovitskiy2021image}, y el Swin Transformer reemplaza la posición absoluta
por un sesgo que depende de la posición relativa entre parches (Sección~\ref{sec:marco_swin}).

\subsection{El bloque transformer}
\label{sec:marco_bloque_transformer}

Un transformer apila bloques iguales. Cada bloque combina una capa de atención multi-cabeza con un perceptrón
multicapa que se aplica a cada token por separado, y rodea a las dos con una conexión residual como la de la
Ecuación~\ref{eq:marco_residual} y una normalización por capa (LN). El transformer original normaliza después
de sumar la conexión residual~\cite{vaswani2017attention}; el \gls{vit}, antes de cada parte,
\begin{equation}
    \mathbf{z}' = \mathbf{z} + \mathrm{Atención}\bigl(\mathrm{LN}(\mathbf{z})\bigr), \qquad
    \mathbf{z}'' = \mathbf{z}' + \mathrm{MLP}\bigl(\mathrm{LN}(\mathbf{z}')\bigr),
    \label{eq:marco_bloque}
\end{equation}
donde $\mathbf{z}$ es la secuencia de embeddings que entra al bloque~\cite{dosovitskiy2021image}. Casi todos
los parámetros del modelo están en las matrices de proyección de la atención y del perceptrón~\cite{kaplan2020scaling}.

El transformer original, planteado para traducción, combina un codificador, que contextualiza la entrada, y un
decodificador, que genera la salida token a token. En el decodificador, una máscara impide que cada posición
atienda a las posteriores, de modo que un token solo depende de los ya generados. Los modelos siguientes
conservan una sola de las dos partes: los \glspl{vit} usan el codificador, y los modelos de lenguaje
generativos de la Sección~\ref{sec:marco_llm_grandes}, el
decodificador~\cite{vaswani2017attention,dosovitskiy2021image,brown2020language}.

\subsection{Vision Transformer}
\label{sec:marco_vit}

El \gls{vit} lleva a las imágenes el codificador del transformer. Divide la imagen en una rejilla de parches
cuadrados, por ejemplo de $16 \times 16$ píxeles (Figura~\ref{fig:marco_vit_patches}), aplana cada parche y lo proyecta a un vector de dimensión
fija, su embedding. A cada parche se le suma el embedding de posición aprendido que indica su lugar en la
rejilla. Sobre esa secuencia de tokens aplica bloques transformer sin operaciones convolucionales, y decide la
categoría a partir de un token de clase que se agrega al inicio de la secuencia y que, al atravesar los bloques,
reúne la información de todos los parches~\cite{dosovitskiy2021image}.

% Figura 2.1. Del parche al token en el Vision Transformer.
% Se usa en sec:marco_vision, subsección "Vision Transformer y Swin Transformer", después del
% párrafo que define el ViT. Cada parche de la fila superior y su token comparten color, para que se
% lea la correspondencia; el token de clase va en gris porque no viene de ningún parche.

\begin{figure}[H]
    \centering
    \begin{tikzpicture}[x=1mm, y=1mm]

        % Imagen dividida en una rejilla de 4x4 parches; la fila superior va coloreada.
        \draw[figtrazo, thick, fill=figrelleno] (-26, -12) rectangle (0, 12);
        \foreach \c [count=\i from 0] in {figazul, fignaranja, figverde, figpurpura} {
            \fill[\c!35] (-26 + \i*6.5, 6) rectangle (-19.5 + \i*6.5, 12);
        }
        \foreach \x in {-19.5, -13, -6.5} { \draw[figtrazo] (\x, -12) -- (\x, 12); }
        \foreach \y in {-6, 0, 6} { \draw[figtrazo] (-26, \y) -- (0, \y); }
        \draw[figtrazo, thick] (-26, -12) rectangle (0, 12);
        \node[font=\small, text=figtrazo, align=center] at (-13, -18) {Imagen dividida\\en parches};

        % Operación de proyección.
        \node[figtarea, text width=28mm] (proj) at (22, 0)
            {Proyectar cada parche\\a un token};

        % Token de clase y tokens de los parches, con su embedding de posición.
        \node[fignodo, minimum width=7mm, minimum height=6mm, inner sep=1pt, fill=figrelleno]
            (t0) at (50, 6) {\footnotesize[CLS]};
        \node[font=\footnotesize, text=figtrazo] at (50, -3) {$+$};
        \node[font=\footnotesize, text=figtrazo] at (50, -9) {$p_{0}$};
        \foreach \c [count=\i from 1] in {figazul, fignaranja, figverde, figpurpura} {
            \node[fignodo, draw=\c!70!black, fill=\c!35, minimum width=6mm, minimum height=6mm,
                  inner sep=1pt] (t\i) at (50 + \i*9, 6) {\footnotesize $x_{\i}$};
            \node[font=\footnotesize, text=figtrazo] at (50 + \i*9, -3) {$+$};
            \node[font=\footnotesize, text=figtrazo] at (50 + \i*9, -9) {$p_{\i}$};
        }
        \node[font=\small, text=figtrazo] at (95, 6) {$\cdots$};
        \node[font=\small, text=figtrazo, align=center] at (72, -18)
            {Secuencia de tokens con\\embedding de posición por parche};

        % Flechas.
        \draw[figflecha] (0, 0) -- (proj.west);
        \draw[figflecha] (proj.east) -- (43, 0) |- (t0.west);

    \end{tikzpicture}
    \caption{Del parche al token en el Vision Transformer, con el token de clase ([CLS]) al inicio.}
    \label{fig:marco_vit_patches}
\end{figure}


El \gls{vit} tiene mucho menos sesgo inductivo propio de las imágenes que una \gls{cnn}: en la convolución, la
localidad y la equivariancia a traslaciones están en cada capa, mientras que en el \gls{vit} solo el
perceptrón de cada bloque es local y la atención es global, de modo que las relaciones espaciales tiene que
aprenderlas de los datos. Preentrenado solo sobre
ImageNet, rinde por debajo de \glspl{cnn} de tamaño comparable; preentrenado sobre
corpus de entre 14 y 300 millones de imágenes, las supera. Por eso importa con qué pesos parte cada
arquitectura antes del ajuste fino (Sección~\ref{sec:marco_transferencia})~\cite{dosovitskiy2021image}.

\subsection{Swin Transformer}
\label{sec:marco_swin}

En el \gls{vit}, cada token atiende a todos los demás, y el costo crece con el cuadrado del número de parches.
El Swin Transformer limita la atención a ventanas locales de $M \times M$ parches. Para un mapa de
$h \times w$ parches con $C$ canales, el costo de la atención global y el de la atención por ventanas son
\begin{equation}
    \Omega_{\text{global}} = 4hwC^2 + 2(hw)^2C, \qquad \Omega_{\text{ventanas}} = 4hwC^2 + 2M^2hwC,
    \label{eq:marco_swin_costo}
\end{equation}
de modo que el segundo crece de forma lineal con el número de parches~\cite{liu2021swin}.

Las ventanas fijas no comunican entre sí. Para que la información cruce de una a otra, los bloques alternan una
partición regular con otra desplazada en medio tamaño de ventana, $\lfloor M/2 \rfloor$ parches, cuyas
ventanas abarcan parches de varias ventanas anteriores (Figura~\ref{fig:marco_swin}). Entre etapas, una capa
de fusión concatena cada grupo de $2 \times 2$ parches vecinos y los proyecta a la mitad de canales que suman,
con lo que la resolución se reduce a la mitad y los canales se duplican. El resultado es una pirámide de mapas
de características, con la estructura jerárquica de una \gls{cnn} y la operación de atención de un
transformer. La posición entra como un sesgo relativo que se suma a la compatibilidad entre cada par de parches
de la ventana~\cite{liu2021swin}.

% Figura: ventanas y ventanas desplazadas del Swin Transformer. Adaptada de liu2021swin (Figura 2).
% Se usa en sec:marco_vit, en el párrafo del Swin.

\begin{figure}[H]
    \centering
    \begin{tikzpicture}[x=1mm, y=1mm]

        % Capa l: cuatro ventanas de 4x4 parches.
        \foreach \c/\x/\y in {figazul/0/16, fignaranja/16/16, figverde/0/0, figpurpura/16/0} {
            \fill[\c!22] (\x, \y) rectangle (\x + 16, \y + 16);
        }
        \draw[figtrazo!40] (0, 0) grid[step=4] (32, 32);
        \draw[figtrazo, very thick] (0, 0) rectangle (32, 32);
        \draw[figtrazo, very thick] (16, 0) -- (16, 32) (0, 16) -- (32, 16);
        \node[font=\footnotesize, text=figtrazo] at (16, -5) {capa $l$: ventanas regulares};

        % Capa l+1: ventanas desplazadas dos parches.
        \begin{scope}[xshift=50mm]
            \foreach \c/\x/\y/\w/\h in {
                figazul/0/24/8/8, fignaranja/8/24/16/8, figverde/24/24/8/8,
                figpurpura/0/8/8/16, figceleste/8/8/16/16, figbermellon/24/8/8/16,
                figazul/0/0/8/8, fignaranja/8/0/16/8, figverde/24/0/8/8} {
                \fill[\c!22] (\x, \y) rectangle (\x + \w, \y + \h);
            }
            \draw[figtrazo!40] (0, 0) grid[step=4] (32, 32);
            \draw[figtrazo, very thick] (0, 0) rectangle (32, 32);
            \draw[figtrazo, very thick] (8, 0) -- (8, 32) (24, 0) -- (24, 32) (0, 8) -- (32, 8) (0, 24) -- (32, 24);
            \node[font=\footnotesize, text=figtrazo] at (16, -5) {capa $l+1$: ventanas desplazadas};
        \end{scope}

        \draw[figflecha] (36, 16) -- (46, 16);

    \end{tikzpicture}
    \caption{Ventanas y ventanas desplazadas del Swin Transformer. Adaptada de \textcite{liu2021swin}.}
    \label{fig:marco_swin}
\end{figure}


\subsection{Redes convolucionales y transformers como línea base}
\label{sec:marco_baseline_visual}

Una línea base (\emph{baseline}) es una arquitectura de referencia, entrenada con protocolo fijo, contra la
que se compara otra: sin ella una exactitud aislada no dice nada, porque es alta o baja según lo que logre la
referencia sobre los mismos datos.

Las \glspl{cnn} y los \glspl{vit} cumplen ese papel para los \glspl{vlm}
(Sección~\ref{sec:marco_vlm}): son las arquitecturas establecidas para clasificar imágenes, con recetas de
entrenamiento maduras y con supervisión completa, es decir, con las etiquetas de la tarea de destino. Un
modelo de otra familia justifica su costo si iguala o supera a esas referencias, o si ofrece algo que ellas
no pueden dar, como producir texto o clasificar sin entrenar sobre la tarea.

La comparación solo es válida si las dos partes usan los mismos datos, particiones y métricas. Puede
enfrentar arquitecturas supervisadas contemporáneas, como ConvNeXt y el
ViT~\cite{dosovitskiy2021image,liu2022convnet}, o familias de modelos, clasificadores supervisados frente a
\glspl{vlm}, y esta segunda dice si el lenguaje aporta algo a la clasificación. Ninguno de estos
clasificadores produce texto. Para eso hace falta un modelo que genere secuencias de palabras, y los modelos de
lenguaje actuales se construyen sobre el decodificador del transformer.
